\section{产品窗口案例复盘：从惊叹到默认工具}

前一章给出挖矿比喻，本章把它拆成产品生命周期，回答“为什么 Magic Moment 很快不够用”这个问题。AI 产品常见路径是：新模型能力出现，创业者包装出第一个惊叹体验，用户快速传播，模型公司和竞品跟进，最后这个能力变成默认功能。只有在传播之后继续沉淀工作流、数据和品牌，产品才不会被全家桶吸收。

\begin{knowledgebox}{Magic Moment 生命周期}
\begin{tabularx}{\textwidth}{p{0.18\textwidth}p{0.34\textwidth}X}
\toprule
阶段 & 发生什么 & 关键动作 \\
\midrule
发现能力 & 新模型能力刚刚可用 & 快速找到最能展示价值的任务。 \\
包装体验 & 做成用户能感知的 Magic Moment & 降低使用门槛，强化传播点。 \\
传播爆发 & 社交媒体和口碑扩散 & 抢默认心智和早期品牌。 \\
竞品跟进 & 模型公司或同类产品复制 & 补工作流、数据和团队协作。 \\
默认化 & 能力变成平台基础功能 & 转向更深场景或成为平台组件。 \\
\bottomrule
\end{tabularx}
\end{knowledgebox}

\begin{importantbox}{从惊叹到默认工具的分水岭}
惊叹来自“第一次看到它能做”；默认工具来自“每天都离不开它”。AI 产品的难点，是在惊叹窗口结束之前完成这次迁移。
\end{importantbox}

\begin{warningbox}{Magic Moment 的副作用}
Magic Moment 会让团队误以为产品已经成立，但用户惊叹不等于长期留存。真正危险的是，团队继续围绕演示效果优化，而不是围绕高频任务、协作流程和付费理由优化。
\end{warningbox}

\subsection{本章小结}

AI 产品的生命周期越来越短，创业者必须把产品从惊叹体验推进到默认工具。否则，Magic Moment 只会成为模型公司全家桶里的一个功能。

